% page 151 of the facsimile (image p-150 of the scan)
% block 170.2 x 250.6 mm ; left margin 31.8 mm
\pgc{10.160mm}{0.000mm}{
\l{}
\l{}
\l{}
\l{\cel{53}141}
\l{}
\l{equiseto. (bot.) Planto kriptogama, di qua la stipo esas rugoza.}
\l{\cel{3}L.\cel{1}\sou{equisetum}. - EL.}
\l{}
\l{equitato. Yusteso naturala, segun la komuna raciono, quan onu}
\l{\cel{3}opozas a la yusteso legala, a la legaro civila. - EFIS.}
\l{}
\l{\sou{equivalar}. (\sou{trans.}) Esar sama-valora (kam). - DEFIRS.}
\l{}
\l{\sou{-er}-. Sufixo qua signifikas \textquotedbl{}qua su okupas ofte pri\textquotedbl{} (ma sen}
\l{\cel{3}agar lo profesionale).}
\l{}
\l{ero. (\sou{kronologio}). I. Epoko konvencionita, de kande onu komencas}
\l{\cel{3}kontar la yari di naciono. - II. Epoko remarkinda, de qua}
\l{\cel{3}komencas koz-ordino nova. - DEFIRS.}
\l{}
\l{\sou{eratika}. (\sou{geol.}) Stranjera en la loko ube onu observas lu, e}
\l{\cel{3}qua konjekteble adportesis ibe de fore (kom ex. grosa, enorma}
\l{\cel{3}fragmenti de roki, e c.) - DEFIS.}
\l{}
\l{\sou{erco}. Substanco minerala qua kontenas un o plura metali. - D.}
\l{}
\l{\sou{erektar}. (\sou{trans.}) Stacigar - per equilibrigar vertikale).-DEFIS.}
\l{}
\l{\sou{ergatoplasmo}. (\sou{biol.}) Parto di la plasmo qua efektigas la laboro}
\l{\cel{3}celulifala - opoze a la trofoplasmo, di qua la rolo esas}
\l{\cel{3}nutrala.}
\l{}
\l{\sou{ergoto}. (\sou{bot.}) Produkturo vejetala parazita, venanta forme di}
\l{\cel{3}mikra pinto, di sporno, sur la spiki di ula graminei, e qua esas}
\l{\cel{3}venenoza. - EFIS.}
\l{}
\l{\sou{ergotino}. (\sou{kemio}). Bazo extraktita de la sekal-ergoto, ed uzata}
\l{\cel{3}medicinale kontre la hemoragii. - DEF.}
\l{}
\l{\sou{ergotismo}. Venenago per la sekalo ergotoza.- DEF.}
\l{}
\l{\sou{-eri-}.\cel{2}Sufixo qua signifikas \textquotedbl{}butiko, butikego en qua onu ofras}
\l{\cel{3}kom varo...\textquotedbl{}, \textquotedbl{}establisuro ube onu (agas to o co).\textquotedbl{}}
\l{}
\l{\sou{eriko}. (\sou{bot.)}\cel{2}Arbusto de la familio \textquotedbl{}erikacei\textquotedbl{}, qua kreskas en}
\l{\cel{3}la suli sabloza. - L.\cel{1}\sou{erica}. - DIL.}
\l{}
\l{\sou{eringo}. (\sou{bot.}) Genero de planto umbelifera, pikantoza.}
\l{\cel{6}\sou{eryngium}.}
\l{}
\l{\sou{eritroblasto}. (\sou{histol.})\cel{2}Globulo reda prima, kun nukleo ed hemo-}
\l{\cel{3}globino, quan onu renkontras dum la periodi di regenero di la}
\l{\cel{3}anemii grava, en la kakexio, e mem quik pos la nasko.}
\l{}
\l{\sou{eritrocito}. (\sou{histol.}) Globulo reda nukleoza, di qua la divideso}
\l{\cel{3}produktas la globuli reda sen-nuklea (hematii).}
\l{}
\l{\sou{eritrofila}. (\sou{biol.}) Qua prizas la koloro reda.}
}
